\documentclass[11pt]{article}

\usepackage[margin=1in]{geometry}
\usepackage[T1]{fontenc}
\usepackage{lmodern}
\usepackage{microtype}
\usepackage{hyperref}
\usepackage{longtable}
\usepackage{booktabs}
\usepackage{array}
\usepackage{enumitem}
\usepackage{xcolor}
\usepackage{amsmath}
\usepackage{amssymb}

\hypersetup{
  colorlinks=true,
  linkcolor=blue!55!black,
  citecolor=blue!55!black,
  urlcolor=blue!55!black,
  pdftitle={Meld Protocol Specification},
  pdfauthor={Meld Project}
}

\newcommand{\meld}{Meld}
\newcommand{\wireversion}{1}
\newcommand{\code}[1]{\ifmmode\text{\texttt{#1}}\else\texttt{#1}\fi}
\newcommand{\field}[1]{\code{#1}}
\newcommand{\usec}{\ensuremath{\mu\mathrm{s}}}
\newcommand{\gf}{\ensuremath{\mathrm{GF}(2^8)}}

\title{\meld{} Protocol Specification\\
\large Working Specification for Wire Version \wireversion}
\author{Meld Project}
\date{June 27, 2026}

\begin{document}
\maketitle

\begin{abstract}
\meld{} is a low-latency coded media transport for live video. It sends
systematic source symbols plus random-linear repair symbols so that a receiver can
recover erasures before a named-packet retransmission can complete a round trip.
The protocol is designed as one deployable adaptive profile, \code{meld-auto},
rather than a collection of static profiles. This document specifies the wire
format, recovery model, feedback loop, media descriptor, security placement, and
implementation requirements for wire version \wireversion.
\end{abstract}

\tableofcontents
\newpage

\section{Status of This Document}

This is a working specification for the current Meld prototype. It is intended to
be precise enough for independent implementation review, internal conformance
tests, and future publication as a formal specification. It is not an Internet
Standards Track document.

Wire version \wireversion{} is the first explicitly versioned Meld wire format.
Earlier prototype headers were not shipped and have no compatibility obligation.
Base-layout changes require a future version nibble. Additive fields are appended
under the extension policy in Section~\ref{sec:versioning}.

\section{Terminology and Conventions}

The key words ``MUST'', ``MUST NOT'', ``REQUIRED'', ``SHALL'', ``SHALL NOT'',
``SHOULD'', ``SHOULD NOT'', ``RECOMMENDED'', ``NOT RECOMMENDED'', ``MAY'', and
``OPTIONAL'' in this document are to be interpreted as described in BCP 14
\cite{rfc2119,rfc8174} when, and only when, they appear in all capitals.

All multi-byte integer fields on the wire are big-endian. Unless explicitly
stated otherwise, clock values and deadlines are signed 64-bit microsecond
timestamps in the sender's clock domain. Fixed-point loss and ECN fractions are
reported as unsigned 16-bit values in parts per 65535.

The following terms are used throughout:

\begin{description}[leftmargin=2.3cm,style=nextline]
\item[Source symbol]
  One fixed-size application chunk after any encryption transform and zero
  padding to \field{SymbolSize}.
\item[Systematic symbol]
  A wire symbol carrying one source symbol verbatim.
\item[Repair symbol]
  A \gf{} linear combination of source symbols.
\item[Sparse repair]
  A repair symbol over an explicit set of source identifiers rather than a
  contiguous coding window.
\item[Flow]
  One logical media stream, identified by a 32-bit \field{Flow} value.
\item[Deadline]
  The latest clock time at which a source symbol can still contribute to
  delivery under the playout budget.
\item[Dependency island]
  A media dependency neighborhood whose loss can damage multiple displayed
  frames, such as a random-access anchor, early reference chain, parameter
  material, or recovery-refresh slice interval.
\end{description}

\section{Design Goals}

\meld{} is optimized for media completion under a playout deadline. The protocol
does not try to beat SRT or RIST at every buffer depth. Its current measured
frontier is iid or tail-erasure loss at tight latency budgets, especially when
the playout budget is less than one RTT. In that region, coded repair can arrive
before ARQ can return a retransmission.

The design goals are:

\begin{enumerate}
\item Recover erased media chunks without waiting for a named-packet ARQ round
      trip when the latency budget is tighter than RTT.
\item Keep one adaptive deployable behavior, \code{meld-auto}; expose latency and
      bitrate budgets, not multiple recovery profiles.
\item Keep source symbols mostly FIFO. Adapt repair timing and rate first.
\item Keep proactive repair inside the configured rate budget unless explicitly
      configured otherwise by an implementation.
\item Use codec-blind transport mechanics, but accept generic media descriptors
      so the transport can estimate decode damage.
\item Prefer conservative behavior when channel classification is uncertain.
\end{enumerate}

SRT and RIST are useful reference points because they are widely deployed
low-latency contribution transports. They primarily recover by retransmission
within a configured latency window, while Meld makes proactive coded recovery the
default low-latency mechanism and uses feedback to tune future repair
\cite{srt-rfc,rist-tr061,rist-tr062}.

\section{Layering Model}

\meld{} has three layers:

\begin{enumerate}
\item The application or media layer writes fixed-size media chunks. It MAY attach
      a generic frame descriptor describing priority, access-unit boundaries,
      references, random-access points, recovery-refresh markers, temporal layer,
      and discardability.
\item The sans-I/O flow core emits source and repair symbols, tracks deadlines,
      estimates loss and burstiness, decodes from any sufficient rank, and emits
      feedback reports. The core does not open sockets and does not read wall
      clocks directly.
\item The session host owns UDP or a caller-provided datagram substrate, pacing,
      timers, optional encryption, DPLPMTUD probes, and goroutines.
\end{enumerate}

The wire waist between implementations is the \field{Symbol} and \field{Feedback}
encoding defined in Section~\ref{sec:wire-format}.

\section{Transport Model}

\meld{} runs over an unreliable datagram substrate. The substrate MAY drop,
duplicate, reorder, or ECN-mark datagrams. The receiver MUST treat corruption as
erasure; corrupted authenticated payloads fail the AEAD check when encryption is
enabled.

Peers are configured with the same \field{SymbolSize}. A coded symbol payload
MUST be exactly \field{SymbolSize} bytes after decoding the Meld header. A
receiver MUST reject a symbol whose payload length differs from its configured
\field{SymbolSize}. Application chunks shorter than \field{SymbolSize} are
zero-padded before coding. Applications that need exact media lengths MUST carry
that length in their media/container format or framing layer.

\section{Protocol Overview}

The sender assigns each source symbol a monotonically increasing 32-bit
\field{SrcIndex}. It emits a systematic symbol immediately, inserts the source
symbol into the live coding window, and emits repair symbols according to its
adaptive controller. Repair symbols are independent equations over the live
source set.

The receiver inserts systematic and repair symbols into an incremental decoder.
Whenever the decoder can recover the next in-order source symbol before its
deadline, the receiver delivers it. If a missing symbol reaches its deadline, the
receiver skips it and advances the delivery cursor so later decoded symbols can
be delivered.

Feedback reports are cumulative state, not events. A lost feedback report is
superseded by the next report. Reports carry delivery progress, rank deficit,
pre-recovery channel loss, burstiness, optional multipath statistics, and optional
media decodability counters.

\section{Versioning and Extension Policy}
\label{sec:versioning}

Every Meld datagram begins with one byte:

\[
  \field{byte0} = (\field{Version} \ll 4) \;|\; \field{Type}
\]

For this specification, \field{Version} MUST be \wireversion. The high nibble
carries the version and the low nibble carries the message type. A decoder MUST
check the version before interpreting the type. If the version is not understood,
the decoder MUST reject the datagram as an unsupported version rather than
attempting to parse it as a different type.

Additive optional fields are appended to the tail of an existing message and
guarded by either a symbol flags bit or a length check. Moving or resizing an
existing base field is an incompatible base-layout change and requires a future
version nibble.

\section{Wire Format}
\label{sec:wire-format}

\subsection{Message Types}

\begin{longtable}{>{\raggedright\arraybackslash}p{0.15\linewidth}
                  >{\raggedright\arraybackslash}p{0.29\linewidth}
                  >{\raggedright\arraybackslash}p{0.45\linewidth}}
\toprule
Type & Message & Description \\
\midrule
\endhead
\code{0x1} & Systematic symbol & One source symbol, verbatim. \\
\code{0x2} & Repair symbol & Random linear combination over a contiguous window. \\
\code{0x3} & Feedback report & Receiver-to-sender cumulative recovery state. \\
\code{0x4} & Clock probe & Receiver-to-sender probe carrying \field{T0}. \\
\code{0x5} & Clock echo & Sender-to-receiver echo carrying \field{T0}, \field{T1}, and \field{T2}. \\
\code{0x6} & Handshake message 1 & Initiator-to-responder encrypted-session handshake payload. \\
\code{0x7} & Handshake message 2 & Responder-to-initiator encrypted-session handshake payload. \\
\code{0x8} & Cookie reply & Responder-to-initiator anti-amplification cookie payload. \\
\code{0x9} & MTU probe & Sender-to-receiver padded DPLPMTUD probe. \\
\code{0xA} & MTU probe acknowledgement & Receiver-to-sender probe acknowledgement. \\
\code{0xB} & Sparse repair symbol & Random linear combination over an explicit source-id set. \\
\bottomrule
\end{longtable}

\subsection{Symbol Header}

Systematic, repair, and sparse repair messages share a 30-byte base header. The
base header is followed by optional extensions, any sparse repair source-id list,
and finally the coded payload.

\begin{longtable}{>{\raggedleft\arraybackslash}p{0.08\linewidth}
                  >{\raggedleft\arraybackslash}p{0.08\linewidth}
                  >{\raggedright\arraybackslash}p{0.22\linewidth}
                  >{\raggedright\arraybackslash}p{0.50\linewidth}}
\toprule
Offset & Size & Field & Meaning \\
\midrule
\endhead
0 & 1 & \field{ver|type} & Version nibble plus symbol type \code{0x1}, \code{0x2}, or \code{0xB}. \\
1 & 4 & \field{Flow} & Flow identifier. \\
5 & 2 & \field{Epoch} & Flow generation and encryption key epoch. \\
7 & 1 & \field{PathID} & Host-stamped path identifier; zero on single-path flows. \\
8 & 4 & \field{WindowBase} & Low edge of the contiguous coding window. \\
12 & 4 & \field{SrcIndex} & Systematic: source id. Repair: repair counter or key-related sequence. \\
16 & 2 & \field{N} & Repair width or sparse-id count. \\
18 & 2 & \field{RepairKey} & Seed for coefficient generation. \\
20 & 1 & \field{Priority} & Protection tier; zero is most disposable. \\
21 & 8 & \field{Deadline} & Decode-by time in sender-clock microseconds. \\
29 & 1 & \field{Flags} & Bit 0: send timestamp extension. Bit 1: frame descriptor extension. \\
\bottomrule
\end{longtable}

If bit 0 of \field{Flags} is set, an 8-byte signed \field{SendTimestamp} follows
the base header. If bit 1 is set, a frame descriptor follows the send timestamp
extension if present, otherwise it follows the base header directly. Extensions
MUST appear in this fixed order.

The frame descriptor extension layout is:

\begin{longtable}{>{\raggedleft\arraybackslash}p{0.10\linewidth}
                  >{\raggedleft\arraybackslash}p{0.10\linewidth}
                  >{\raggedright\arraybackslash}p{0.22\linewidth}
                  >{\raggedright\arraybackslash}p{0.46\linewidth}}
\toprule
Offset & Size & Field & Meaning \\
\midrule
\endhead
0 & 4 & \field{FrameStart} & Source id of the first chunk in the access unit. \\
4 & 2 & \field{FrameLen} & Number of chunks in the access unit. \\
6 & 1 & \field{DescFlags} & Bit 0 RAP, bit 1 discardable, bit 2 non-picture, bit 3 recovery-refresh. \\
7 & 1 & \field{nRefs} & Number of referenced access units. \\
8 & 4 each & \field{RefStart[]} & Source id of each referenced access unit's first chunk. \\
\bottomrule
\end{longtable}

A sender SHOULD attach the frame descriptor only to systematic symbols. Coding
recovers payload bytes, not headers, so a recovered source symbol does not carry
the descriptor extension. Receivers compute frame decodability from directly
observed descriptors and source-id ranges.

\subsection{Systematic Symbols}

A systematic symbol uses type \code{0x1}. It carries one source symbol verbatim.
\field{SrcIndex} is the source identifier. \field{N} SHOULD describe the width of
the current generation or be one on the sliding path. \field{RepairKey} is unused
by the systematic equation and SHOULD be zero.

\subsection{Contiguous Repair Symbols}

A contiguous repair symbol uses type \code{0x2}. \field{WindowBase} and \field{N}
identify the source ids covered by the repair equation:

\[
  [\field{WindowBase}, \field{WindowBase} + \field{N})
\]

\field{RepairKey} seeds the deterministic coefficient generator. The resulting
payload is the byte-wise \gf{} linear combination of the covered source-symbol
payloads.

\subsection{Sparse Repair Symbols}

A sparse repair symbol uses type \code{0xB}. It carries a repair equation over an
explicit list of source ids. After the optional symbol extensions and before the
payload, the datagram carries \field{N} unsigned 32-bit source identifiers in
coefficient order. \field{N} MUST be greater than zero and MUST NOT exceed 64.

Sparse repair is intended for protected dependency neighborhoods where a
contiguous window would spend repair pressure on unrelated disposable symbols.
It is still part of \code{meld-auto}; it is not a separate profile.

\subsection{Feedback Report}

A feedback report uses type \code{0x3}. The base report is 53 bytes. Additional
tail groups are length-gated.

\begin{longtable}{>{\raggedleft\arraybackslash}p{0.08\linewidth}
                  >{\raggedleft\arraybackslash}p{0.08\linewidth}
                  >{\raggedright\arraybackslash}p{0.24\linewidth}
                  >{\raggedright\arraybackslash}p{0.48\linewidth}}
\toprule
Offset & Size & Field & Meaning \\
\midrule
\endhead
0 & 1 & \field{ver|type} & Version nibble plus \code{0x3}. \\
1 & 4 & \field{Flow} & Flow identifier. \\
5 & 2 & \field{Epoch} & Flow generation, mirroring symbol epoch. \\
7 & 4 & \field{DecodedLowEdge} & All source ids below this value are delivered or skipped. \\
11 & 4 & \field{HighestSeen} & Highest observed source-id edge. \\
15 & 2 & \field{Deficit} & Extra independent symbols needed at the cursor. \\
17 & 2 & \field{EcnCE} & CE-marked fraction this interval, parts per 65535. \\
19 & 2 & \field{LossRate} & Smoothed pre-recovery erasure estimate, parts per 65535. \\
21 & 32 & \field{Deficits[32]} & Rank deficits for the next 32 generations or windows, saturating at 255. \\
53 & 2 & \field{CongestionLoss} & Pre-recovery wire-loss count since the last report. \\
55 & 2 & \field{Burstiness} & Mean loss-run length in Q8; 256 denotes iid. \\
57 & 1 & \field{nPaths} & Number of paths in the multipath tail; zero means absent or single path. \\
\bottomrule
\end{longtable}

When \field{nPaths} is non-zero, it MUST be at most 8. It is followed by
\field{PathLoss[nPaths]} as uint16 fractions and \field{SlotDist[nPaths+1]} as a
histogram of aligned-slot erasure counts. \field{SlotDist[j]} is the fraction of
slots in which exactly \(j\) of \(n\) paths erased their symbol.

After the per-path section, a 16-byte media-damage tail MAY be present. It
contains four uint32 counters: \field{Frames}, \field{DecodableFrames},
\field{Keyframes}, and \field{Decodable\allowbreak{}Keyframes}. A sender uses
these counters to distinguish ordinary packet loss from dependency damage.

Feedback is cumulative. A sender MUST treat absent tail fields as zero or unknown,
not as protocol failure.

\subsection{Clock Probe and Echo}

A clock probe uses type \code{0x4} and has length 9 bytes:
\field{ver|type} followed by \field{T0} as a signed 64-bit timestamp.

A clock echo uses type \code{0x5} and has length 25 bytes:
\field{ver|type}, \field{T0}, \field{T1}, and \field{T2}. With receiver receive
time \field{T3}, the receiver estimates offset and RTT using:

\[
  \mathrm{offset} = \frac{(T1 - T0) + (T2 - T3)}{2}
\]
\[
  \mathrm{rtt} = (T3 - T0) - (T2 - T1)
\]

\subsection{DPLPMTUD Probes}

An MTU probe uses type \code{0x9}. It carries a 4-byte nonce and zero padding so
the whole datagram has the candidate size under test. An MTU probe acknowledgement
uses type \code{0xA} and carries the nonce plus the observed datagram size as a
uint16. This follows the datagram packetization-layer PMTUD model of RFC 8899
\cite{rfc8899}.

\subsection{Handshake Payload Framing}

Handshake messages are framed by the Meld version/type byte followed by a
suite-specific crypto payload. The wire package treats the payload as opaque; the
current encrypted-session suite pins the following payload lengths:

\begin{longtable}{>{\raggedright\arraybackslash}p{0.12\linewidth}
                  >{\raggedright\arraybackslash}p{0.18\linewidth}
                  >{\raggedleft\arraybackslash}p{0.16\linewidth}
                  >{\raggedright\arraybackslash}p{0.42\linewidth}}
\toprule
Type & Direction & Payload bytes & Payload layout \\
\midrule
\endhead
\code{0x6} & Init to resp & 1248 & X25519 public key (32), ML-KEM-768 encapsulation key (1184), \field{mac1} (16), \field{mac2} (16). \\
\code{0x7} & Resp to init & 1152 & X25519 public key (32), ML-KEM-768 ciphertext (1088), confirmation tag (16), \field{mac1} (16). \\
\code{0x8} & Resp to init & 56 & Cookie nonce (24), encrypted cookie body (32). \\
\bottomrule
\end{longtable}

\section{Coding Semantics}

\meld{} uses random linear network coding over \gf{}. A source symbol is a unit
vector equation. A repair symbol is an equation whose coefficient vector is
derived deterministically from \field{RepairKey} and the covered source ids.

For contiguous repair, the covered ids are the consecutive ids named by
\field{WindowBase} and \field{N}. For sparse repair, the covered ids are the
explicit \field{SparseIDs} list. A receiver MUST treat a repair symbol as useful
only if its equation is innovative relative to the receiver's current linear
system.

The default low-latency path is sliding-window coding. The implementation MAY
retain a generation-mode fallback, but peers MUST interoperate through the same
symbol and feedback waist. A receiver follows the sender-stamped
\field{WindowBase} and \field{N}; it MUST NOT assume a fixed generation width
when stamped widths vary.

\section{Sender Requirements}

A conforming sender:

\begin{enumerate}
\item MUST assign source ids monotonically within an epoch.
\item MUST emit source symbols in source-id order except when an implementation
      intentionally sheds a discardable media layer before source assignment.
\item MUST NOT depend on full emitted-stream lookahead for sender-realizable
      placement. Repair timing decisions are limited to symbols already created
      or safely releasable before deadline risk.
\item SHOULD size proactive repair from measured erasure rate, burstiness, and a
      target decode-failure probability.
\item SHOULD send reactive top-up repair when feedback reports a live rank
      deficit and the repair can still arrive before the relevant deadline.
\item MUST cap or throttle repair when a configured aggregate rate budget would
      otherwise be exceeded, unless an operator has explicitly disabled that cap.
\item SHOULD preserve source FIFO under congestion and reduce repair before
      delaying non-discardable source past deadline.
\end{enumerate}

When uncertainty increases, for example after missing feedback, high reorder, or
unstable loss estimates, the sender SHOULD bias conservative: keep source FIFO,
avoid speculative placement, and rely on feedback-driven repair that can still
land in time.

\section{Receiver Requirements}

A conforming receiver:

\begin{enumerate}
\item MUST reject unknown wire versions.
\item MUST reject symbols for the wrong flow.
\item MUST reject coded payloads whose length differs from configured
      \field{SymbolSize}.
\item MUST bound decode state so a forged source id or window cannot force
      unbounded allocation or cursor walking.
\item MUST deliver source symbols in source-id order.
\item MUST skip a missing source symbol when its deadline expires, so later
      recovered symbols are not blocked indefinitely.
\item SHOULD report feedback periodically and after meaningful recovery-state
      changes.
\end{enumerate}

Receivers SHOULD count pre-recovery wire loss from systematic source-id gaps
before successful decoding hides those gaps. This preserves an honest congestion
and repair-sizing signal.

\section{Adaptive Control Loop}

\code{meld-auto} is the only deployable profile. It classifies the current
frontier online using:

\begin{itemize}
\item RTT and one-way delay variation from clock exchange and timestamps,
\item playout budget and slack after RTT,
\item pre-recovery loss rate,
\item loss-run burstiness,
\item reorder behavior,
\item optional multipath erasure correlation, and
\item optional media decodability damage.
\end{itemize}

The sender uses those signals to choose repair aggressiveness, sparse repair
eligibility, reactive repair timing, rate-budget throttling, temporal-layer
shedding if enabled, and an optional encoder recovery-cadence request. The
control loop is continuous; changing network conditions should move the sender
between regimes without requiring a named user profile.

\section{Media Descriptor}

The media descriptor is generic and codec-blind. It carries enough structure for
the receiver to compute decode damage without parsing AVC, HEVC, AV1, JPEG XS, or
container-specific syntax.

The descriptor fields are:

\begin{description}[leftmargin=3.0cm,style=nextline]
\item[\field{Priority}]
  Protection tier. Higher tiers receive tighter decode-failure targets.
\item[\field{FrameStart}]
  First source id of the access unit.
\item[\field{FrameLen}]
  Number of chunks in the access unit.
\item[\field{RefStart[]}]
  First source ids of dependency access units.
\item[RAP]
  Random-access point or equivalent recovery anchor.
\item[Recovery-refresh]
  Reference slice or access unit participating in a bounded intra-refresh or
  recovery interval.
\item[Discardable]
  Unit that no later unit depends on.
\item[Non-picture]
  Parameter, metadata, or other non-displayed coded material.
\end{description}

The descriptor does not mandate a codec structure. A media shaper maps codec or
container facts into this generic form.

\section{Encoder Recovery-Cadence Actuator}

\meld{} MAY expose an advisory encoder-control output. When feedback indicates
that long bursts are damaging dependency islands, the sender MAY request a
bounded maximum recovery interval from an attached encoder. An encoder can satisfy
that request with bounded intra-refresh, recovery-point signaling, keyframes, or
another codec-native mechanism.

This actuator is advisory. It MUST NOT create multiple deployed Meld profiles. If
the encoder cannot comply, Meld continues with the same transport loop and
configured repair budget.

\section{Multipath}

The \field{PathID} field identifies the path on which a symbol was sent. In a
single-path session it is zero. In a multipath session, systematic and repair
symbols are spread across paths and decoded from their union. The receiver reports
per-path marginal loss and an aligned-slot erasure-count histogram. The sender
uses this histogram to size repair against the joint erasure tail rather than
assuming independent paths.

Path count and path layout are deployment configuration, not wire-negotiated in
version \wireversion. A receiver SHOULD stop reporting multipath correlation
statistics if observed \field{PathID} stamps show that peers disagree about the
path layout.

\section{Pacing and Deadlines}

The session host SHOULD pace datagrams to the configured aggregate rate budget.
Pacing MUST be release-constrained: it can only send datagrams that already exist
and are safe to hold until the release time. A sender MUST NOT model pacer
reordering as if it had full future-stream lookahead.

\field{Deadline} is per symbol. Repair stamped with a deadline is useful only if
it can arrive before the source symbols it protects expire. A receiver MAY clamp
unreasonable peer-stamped deadlines to a sane window around local time to protect
against forged timestamps.

\section{Security}

Encryption is optional. When enabled, it lives in the session host and the flow
core remains crypto-blind.

\subsection{Handshake}

The encrypted session handshake uses an Argon2id-stretched passphrase as a PSK,
an NNpsk0-style Noise pattern, X25519, ML-KEM-768, HKDF-SHA256, and
ChaCha20-Poly1305. Handshake message bytes are carried as opaque payloads behind
message types \code{0x6}, \code{0x7}, and \code{0x8}. The cookie reply is an
anti-amplification mechanism used under load.

\subsection{Encrypt-then-code}

The host AEAD-seals each source symbol before the coder sees it. The coder then
treats ciphertext plus tag as opaque payload bytes. Repair symbols are linear
combinations of ciphertext bytes. The receiver solves the linear system, recovers
source ciphertexts, and then authenticates and decrypts each source symbol.

This preserves end-to-end confidentiality and integrity while allowing a relay to
route or recode without plaintext access. The symbol header remains visible. In
encrypted mode, source-symbol identity fields are authenticated as AEAD associated
data.

\subsection{Control Authentication}

On encrypted flows, feedback, clock, and MTU control datagrams are authenticated
with a trailer carrying an 8-byte sequence number and a 16-byte HMAC-SHA256 tag.
Receivers drop replayed or stale control datagrams using a bounded replay window.

\subsection{Limitations}

Version \wireversion{} does not provide anonymity, metadata confidentiality,
group keying, or per-hop pollution detection for untrusted recoding relays.
Cleartext headers reveal flow id, window structure, priority, deadlines, frame
boundaries, and dependency hints.

\section{Congestion and Rate Budget}

\meld{} distinguishes pre-recovery wire loss from post-recovery media delivery.
Coding MUST NOT hide congestion signals from the sender. The receiver reports
pre-recovery loss and ECN CE fraction; the sender uses those signals to size
repair and, when congestion control is enabled, to adjust its budget.

The sender SHOULD keep proactive repair inside \field{MaxBitrate}. If the budget
cannot carry both source and repair, source and non-discardable base media take
priority over repair. Optional media shedding is limited to discardable temporal
layers and occurs before source-id assignment so it does not create transport
holes.

\section{Benchmarking Requirements}

Meld performance claims SHOULD be evaluated at the glass-to-glass level:
decoded frames, decodable frame percentage, decodable keyframe percentage, same
source bytes and packets across transports, RTT, playout budget relative to RTT,
loss model, and oracle rows.

The current curated frontier evidence is non-normative but important:

\begin{itemize}
\item At 10\% iid loss with a 0.75x RTT playout budget across 50, 100, and
      200 ms RTT cells, \code{meld-auto} reached 144.0 decoded frames while the
      best ARQ row, RIST, reached 121.1 frames, a +22.9 frame delta.
\item At 5\% iid loss with the same 0.75x RTT budget, \code{meld-auto} reached
      143.8--144.0 frames and held a roughly +10 frame advantage.
\item At 1x RTT and 1.5x RTT budgets, ARQ catches up and the Meld advantage
      compresses.
\item Burst48 and Burst96 experiments did not produce a stable deployable Meld
      target under the current source structure and repair budget.
\end{itemize}

New optimizations SHOULD be judged by movement of the \code{meld-auto} macro
frontier against SRT, RIST, and oracle rows rather than isolated packet-level
improvements.

\section{Interoperability Checklist}

An independent implementation of wire version \wireversion{} should verify:

\begin{enumerate}
\item Version/type nibble parsing rejects unknown versions before unknown types.
\item Symbol, feedback, clock, handshake, and MTU message types round-trip.
\item Systematic and repair payloads are exactly \field{SymbolSize}.
\item Contiguous and sparse repair equations use identical coefficient generation
      for the same source ids and \field{RepairKey}.
\item Feedback tails are decoded only when present and absent tails default to
      zero or unknown.
\item Frame descriptors preserve source-id ranges and reference starts.
\item Deadlines skip missing source ids without blocking later recovered ids.
\item Encrypted flows seal before coding and open after decoding.
\end{enumerate}

\appendix

\section{Relationship to Existing Protocols}

SRT specifies a UDP-based reliable transport with data and control packet types,
latency-window behavior, ACK/NAK feedback, and retransmission-oriented recovery
\cite{srt-rfc}. RIST profiles specify interoperable reliable internet streaming
over RTP/RTCP mechanisms with retransmission and optional features layered by
profile \cite{rist-tr061,rist-tr062}. Meld borrows the discipline of precise
packet formats, cumulative feedback, and deployment-oriented control surfaces,
but uses random-linear coded recovery as the primary low-latency loss recovery
tool.

The Meld coding format follows the same broad lineage as the RLC FEC scheme in
RFC 8681: compact repair identifiers regenerate coding coefficients, and repair
symbols can be consumed as soon as they arrive \cite{rfc8681}. Meld differs by
embedding the coding waist directly into a media transport with deadlines,
media-damage feedback, and an adaptive single-profile controller.

\begin{thebibliography}{99}

\bibitem{rfc2119}
S. Bradner, ``Key words for use in RFCs to Indicate Requirement Levels,''
BCP 14, RFC 2119, March 1997. \url{https://www.rfc-editor.org/rfc/rfc2119}

\bibitem{rfc8174}
B. Leiba, ``Ambiguity of Uppercase vs Lowercase in RFC 2119 Key Words,''
BCP 14, RFC 8174, May 2017. \url{https://www.rfc-editor.org/rfc/rfc8174}

\bibitem{rfc8681}
V. Roca, B. Teibi, and C. Burdinat, ``Sliding Window Random Linear Code
(RLC) Forward Erasure Correction (FEC) Scheme for FECFRAME,'' RFC 8681,
January 2020. \url{https://www.rfc-editor.org/rfc/rfc8681}

\bibitem{rfc8899}
G. Fairhurst, T. Jones, M. Tuexen, I. Ruengeler, and T. Voelker, ``Packetization
Layer Path MTU Discovery for Datagram Transports,'' RFC 8899, September 2020.
\url{https://www.rfc-editor.org/rfc/rfc8899}

\bibitem{srt-rfc}
SRT Alliance, ``The SRT Protocol,'' Internet-Draft style technical
specification. \url{https://haivision.github.io/srt-rfc/draft-sharabayko-srt.html}

\bibitem{rist-tr061}
Video Services Forum, ``Technical Recommendation TR-06-1:2020 Reliable Internet
Stream Transport (RIST) Simple Profile.'' \url{https://vsf.tv/technical-recommendations/}

\bibitem{rist-tr062}
Video Services Forum, ``Technical Recommendation TR-06-2:2024 Reliable Internet
Stream Transport (RIST) Main Profile.'' \url{https://vsf.tv/technical-recommendations/}

\end{thebibliography}

\end{document}
